\paragraph{Replication.} Both simulated-user experiments of \citet{ribeiro2016} replicate. With the authors' algorithms and settings ported to Python~3, LIME is the best of the four explainers in all 16 settings. Its trust F1 is within \ReplTrustMAE{} points of the published values on average, and its faithfulness recall is \ReplRecallGap{} points lower on average. The largest discrepancy, the Parzen baseline's faithfulness, is partly explained by bandwidths tuned to the authors' unavailable pre-tokenised features (our ablation) and otherwise by that input change. The paper is therefore reproducible in substance but not exactly: exact reproduction would need the frozen Python~2 stack and the \texttt{processed\_acl} files, and it is this dependency, not the method, that has decayed.

\paragraph{Generalisation.} On a new text domain (IMDB), every conclusion carries over. On a new modality (tabular clinical data), two do not: LIME is much less faithful to decision trees (\HeartTreeRecall{}\% recall), and a simple Parzen baseline beats it at telling trustworthy from untrustworthy predictions. We trace the second result to what LIME's tabular explanations mean, deviation from the training distribution rather than from a reference value. For tabular data, a practitioner should therefore read LIME weights as ``relative to a typical patient'', not ``what would happen if this value were removed''. With real people the picture is more cautious: raters agreed well and trusted correct predictions more, but the model's displayed confidence explains all of that discrimination, so in our data the explanations themselves did not measurably help people decide when to trust a prediction. An LLM judge reproduced the humans' ranking but accepted every explanation, so it cannot yet replace human evaluation of explanation quality.

\paragraph{Reflections.} (1)~\emph{Reproducibility decays through the environment more than through the paper.} The method and evaluation were described clearly enough to re-implement, but the pinned 2016 stack no longer installs and one input format is no longer linked. Version-pinned containers and released processed data would have avoided both problems. (2)~\emph{Protocols do not transfer between modalities for free.} The paper's notion of ``removing'' a feature is clear for words but not for a blood-pressure reading. Our choice (reset to the median or mode) is defensible, but it directly affects which explainer looks best, which suggests that explanation benchmarks should state their removal semantics explicitly. (3)~\emph{Annotation is harder than it looks.} Raters without clinical training must judge both the model and the medicine, and moderate agreement is expected. In hindsight, a larger set of items with simpler, non-medical content would separate ``trust in the explanation'' from ``trust in one's own medical intuition'' more cleanly. A design where each rater sees the same item both with and without an explanation would isolate the explanation's effect. (4)~\emph{Our earlier pipeline gave a more flattering number.} The project-update run reported a trust F1 of 97.5 on one setting with a simplified trust rule. The faithful port gives 95.5. Small protocol simplifications can bias a replication towards ``success'', which is why we kept the original rules even where they look odd (e.g.\ explaining the true rather than the predicted label in \S5.2).
